% !TeX program = xelatex
\documentclass[11pt,a4paper]{article}
\usepackage[utf8]{inputenc}
\usepackage[T1]{fontenc}
\usepackage[english]{babel}
\usepackage{amsmath,amssymb,amsfonts,mathtools,amsthm}
\usepackage{geometry}
\geometry{left=1.25cm,right=1.25cm,top=1.25cm,bottom=3.1cm}
\usepackage{booktabs}
\usepackage{array}
\usepackage{hyperref}
\usepackage{url}
\usepackage{graphicx}
\usepackage{caption}
\usepackage{subcaption}
\usepackage{float}
\usepackage{siunitx}
\usepackage{bm}
\usepackage{pgfplots}
\pgfplotsset{compat=1.18}
\usepackage{xcolor}
\usepackage{titlesec}
\usepackage{fancyhdr}
\usepackage{microtype}
\usepackage{fontspec}
\usepackage{tcolorbox}
\usepackage{enumitem}
\usepackage{xeCJK}  % 中文支持

% 设置拉丁字体（同原模板）
\setmainfont{Calibri}
\setsansfont{Segoe UI}[
BoldFont = Segoe UI Bold,
Ligatures = TeX
]

% 中文字体
\setCJKmainfont{Microsoft YaHei}
\setCJKsansfont{Microsoft YaHei}
\setCJKmonofont{Microsoft YaHei}

\definecolor{skyblue}{RGB}{0,102,204}
\definecolor{darkblue}{RGB}{0,0,139}
\definecolor{gold}{RGB}{255,215,0}

% YUCT 宏
\newcommand{\Keff}{K_{\mathrm{eff}}}
\newcommand{\kappac}{\kappa_c}
\newcommand{\eps}{\varepsilon}
\newcommand{\betaf}{\beta}
\newcommand{\df}{d_{\!f}}
\newcommand{\Sodd}{S_{\mathrm{odd}}}
\newcommand{\Seven}{S_{\mathrm{even}}}
\newcommand{\Mpl}{M_{\mathrm{Pl}}}
\newcommand{\Lpl}{l_{\mathrm{Pl}}}
\newcommand{\tpl}{t_{\mathrm{Pl}}}
\newcommand{\Lmin}{\ell_{\min}}
\newcommand{\LUV}{\Lambda_{\mathrm{UV}}}

% 定理环境（中文）
\newtheorem{theorem}{定理}[section]
\newtheorem{proposition}[theorem]{命题}
\newtheorem{definition}[theorem]{定义}
\theoremstyle{remark}
\newtheorem{remark}[theorem]{备注}

% 章节格式
\titleformat{\section}{\LARGE\bfseries\color{darkblue}}{\thesection}{1em}{}
\titleformat{\subsection}{\Large\bfseries\color{darkblue}}{\thesubsection}{1em}{}
\titleformat{\subsubsection}{\normalsize\bfseries\color{darkblue}}{\thesubsubsection}{1em}{}

% 页眉页脚（与中文模板一致）
\fancypagestyle{firstpage}{%
	\fancyhf{}%
	\renewcommand{\headrulewidth}{-5pt}%
	\renewcommand{\footrulewidth}{-5pt}%
	\fancyfoot[C]{%
		\begin{minipage}[t]{\textwidth}
			\centering
			\textcolor{gold}{\rule{\textwidth}{1.6pt}}\\[5pt]
			{\small \textsuperscript{1}Yakushev Research, \url{https://yuct.org/}} \hfill
			{\small YPSDC 协议: \url{https://ypsdc.com/}}\\[-2pt]
			\textcolor{gold}{\rule{\textwidth}{1.6pt}}\\[5pt]
			\textbf{雅库舍夫统一协调理论 2026} \hfill \thepage\\[10pt]
		\end{minipage}%
	}%
}

\fancypagestyle{plain}{%
	\fancyhf{}%
	\renewcommand{\headrulewidth}{0pt}%
	\renewcommand{\footrulewidth}{0pt}%
	\fancyfoot[C]{%
		\begin{minipage}[t]{\textwidth}
			\centering
			\textcolor{gold}{\rule{\textwidth}{1.6pt}}\\[4pt]
			\textbf{雅库舍夫统一协调理论 2026} \hfill \thepage\\[4pt]
		\end{minipage}%
	}%
}

\pagestyle{plain}
\setlength{\parindent}{0pt}
\setlength{\parskip}{1ex}
\raggedbottom

\hypersetup{
	colorlinks=true,
	linkcolor=blue,
	citecolor=blue,
	urlcolor=blue,
}

% 自定义页眉（右侧显示“附录 UVD”）
\newcommand{\makeheader}{%
	\noindent
	\begin{minipage}{0.17\textwidth}
		\raggedright
		{\fontspec{Segoe UI}[BoldFont=Segoe UI Bold]\bfseries\color{skyblue}\fontsize{46}{46}\selectfont YUCT}%
	\end{minipage}%
	\hspace{30pt}
	\begin{minipage}{0.32\textwidth}
		\raggedright
		{\sffamily\fontsize{11}{13}\selectfont YAKUSHEV\\ UNIFIED\\ COORDINATION THEORY}%
	\end{minipage}%
	\hfill
	\begin{minipage}{0.38\textwidth}
		\raggedleft
		{\sffamily\bfseries 附录 UVD}\\ 
		{\sffamily 版本 1.0 / 2026年6月}\\ 
		{\sffamily\footnotesize \url{https://doi.org/10.5281/zenodo.18444598}}%
	\end{minipage}
	\par\vspace{5pt}
	\noindent\textcolor{gold}{\rule{\textwidth}{1.6pt}}
	\par\vspace{0pt}
}

\begin{document}
	
	\thispagestyle{firstpage}
	\makeheader
	
	\vspace{0.5cm}
	{\LARGE\bfseries 附录 UVD：YUCT 中紫外线发散问题的解决 \\
		\Large 来自协调形状因子和协调网络物理标度的有限量子场论 \par}
	\vspace{0.3cm}
	
	{\large Alexey V. Yakushev\textsuperscript{1}} \\
	\vspace{0.2cm}
	
	{\color{darkblue}\textbf{\LARGE 摘要}} \par
	{\large
		\noindent
		雅库舍夫统一协调理论（YUCT）提出，量子场论中的紫外线发散通过一个协调形状因子 \(F(q^2) = \exp[-(q^2/\LUV^2)^{2/3}]\) 消除，该因子源于真空的有限信息容量。在最初的表述中，紫外尺度 \(\LUV\) 被等同于普朗克质量，\(\LUV = 3\pi\Mpl \approx 1.15\times 10^{20}\) GeV，导致对希格斯质量产生巨大的有限辐射修正。我们修正了标量质量修正中出现的常数 \(C_0\) 的计算错误：正确值为 \(C_0 = 3\sqrt{\pi}/(8\sqrt{2}) \approx 0.46999\)，取代了先前报告的 \(C_0 \approx 0.626\)。通过这一修正，来自普朗克尺度网络的辐射贡献为 \(\delta m \approx 1.57\times 10^{18}\) GeV，这重新引入了严重的精细调节问题，除非裸质量被精确抵消。我们考察两种逻辑上可行的替代方案：（i）YUCT 质量公式 \(m = \Mpl(2/3)^L\xi_f\) 已经确定了物理质量，而辐射修正被一个有限的反项吸收——这在技术上自洽，但审美上不令人满意；（ii）普朗克尺度下的基本拉格朗日量不包含裸质量项（标度不变性），整个希格斯质量完全由辐射修正产生。在这种第二个情景中，自洽性要求紫外尺度为 \(\LUV \approx 8.96\) TeV，这是 LHC 能够达到的值。然后协调形状因子预测在 TeV 尺度的高能散射截面中会出现可观测的偏差。我们呈现这两种情景、它们的数学论证以及实验后果，目的是让 YUCT 接受严格的实证检验。
	}
	
	\vspace{0.2cm}
	\noindent
	{\color{darkblue}\textbf{关键词:}} YUCT，紫外线发散，协调形状因子，等级问题，希格斯质量，标度不变性，非定域量子场论，LHC 现象学。
	\vspace{0.1cm}
	\noindent
	{\color{darkblue}\textbf{引用:}} Yakushev AV. 附录 UVD：YUCT 中紫外线发散问题的解决. 2026. DOI: \url{https://doi.org/10.5281/zenodo.18444598}（本卷）.
	
	\vspace{0.5cm}
	
	% ============================================================
	\section{引言}
	\label{sec:intro}
	
	雅库舍夫统一协调理论（YUCT）为量子场论（QFT）中的紫外线发散问题提供了一个优雅的解决方案。YUCT 不是引入任意的动量截断或无穷大重整化，而是假设物理真空具有一个有限的用于传输量子信息的“带宽”。该带宽由一个\textbf{协调形状因子} \(F(q^2)\) 量化，该形状因子在特征尺度 \(\LUV\) 以上指数压制虚动量。这个形状因子的数学结构由 YUCT 的普适误差律 \(\varepsilon \propto K_{\mathrm{eff}}^{-\beta}\)（其中 \(\beta = 2/3\)）决定。
	
	在此附录的先前版本中，我们将 \(\LUV\) 等同于普朗克质量，\(\LUV = 3\pi\Mpl \approx 1.15\times 10^{20}\) GeV——该尺度是从 YUCT 的基本代数循环所要求的极小长度 \(\Lmin = \frac{2}{3}\Lpl\) 推导出来的。在这种等同下，所有圈积分都变得明显有限：对数发散被替换为 \(\ln(\LUV/m)\)，而幂次发散完全消失。
	
	然而，对标量质量修正的仔细重算揭示出常数 \(C_0\) 的一个计算错误，该常数决定了有限残余的大小。\(C_0\) 的正确值小于先前报告的值，这减小了辐射修正，但仍将其保持在比观测到的希格斯质量大许多数量级的水平。这迫使我们重新审视物理尺度 \(\LUV\)。
	
	本附录有两个目标。首先，我们详细呈现修正后的数学推导，确保所有表达式准确且可重现。其次，我们探讨修正后公式的物理含义，包括一个激进的提议：协调网络的基本尺度可能根本不是普朗克尺度，而是一个低得多的尺度，可能处于当前对撞机实验可及的范围内。
	
	% ============================================================
	\section{协调形状因子与紫外尺度}
	\label{sec:formfactor}
	
	YUCT 协调形状因子为
	\begin{equation}
		F(q^2) = \exp\!\Bigl[-\Bigl(\frac{q^2}{\LUV^2}\Bigr)^{\!\beta}\,\Bigr], \qquad \beta = \frac{2}{3}.
		\label{eq:F}
	\end{equation}
	所有标准模型的传播子都乘以 \(F(p^2/\LUV^2)\)。尺度 \(\LUV\) 最初是通过康普顿关系 \(\LUV = 2\pi\hbar c/\Lmin = 3\pi\Mpl\) 从极小长度 \(\Lmin = \frac{2}{3}\Lpl\) 推导出来的。
	
	这个形状因子使得所有圈积分在紫外区域绝对收敛。证明很简单：对于圈动量中的任何多项式 \(P(k)\)，乘积 \(P(k) \cdot \prod_i F(k_i^2/\LUV^2)\) 当 \(|k| \to \infty\) 时下降速度快于 \(|k|\) 的任何负幂次，因为指数函数占主导。
	
	% ============================================================
	\section{电子自能：对数发散问题的解决}
	\label{sec:electron}
	
	在常规 QFT 中，QED 的单圈电子自能是对数发散的。引入 YUCT 形状因子后，欧几里得积分变为
	\begin{equation}
		\Sigma_{\mathrm{YUCT}} \propto \int_0^{\infty} \frac{dk_E}{k_E}\; \exp\!\Bigl[-2\Bigl(\frac{k_E^2}{\LUV^2}\Bigr)^{\!2/3}\,\Bigr].
		\label{eq:elec_int}
	\end{equation}
	代换 \(t = (k_E^2/\LUV^2)^{2/3}\) 给出 \(dk_E/k_E = \frac{3}{4} dt/t\)，积分结果为
	\begin{equation}
		I_{\mathrm{log}} = \frac{3}{4} \int_{t_{\min}}^{\infty} \frac{dt}{t}\, e^{-2t} \approx \ln\frac{\LUV}{m_e},
		\label{eq:elec_result}
	\end{equation}
	其中 \(t_{\min} \sim (m_e^2/\LUV^2)^{2/3}\)。对数发散被替换为有限对数 \(\ln(\LUV/m_e) \approx 53.8\)。这个计算是初等的，不需要修正。
	
	% ============================================================
	\section{标量质量修正：常数 \(C_0\) 的修正}
	\label{sec:scalar}
	
	\subsection{积分及其计算}
	
	在 \(\phi^4\) 理论中，经过维克旋转并代入形状因子后，标量场的单圈质量修正为
	\begin{equation}
		\delta m^2_{\mathrm{YUCT}} = \frac{\lambda}{16\pi^2} \int_0^{\infty} \frac{k_E^3\, dk_E}{k_E^2 + m^2}\; F^2\!\Bigl(\frac{k_E^2}{\LUV^2}\Bigr), \qquad F^2(x) = e^{-2x^{2/3}}.
		\label{eq:scalar_YUCT}
	\end{equation}
	
	引入无量纲变量 \(t = k_E^2/\LUV^2\)，并在 \(m \ll \LUV\) 条件下展开，得到主导项
	\begin{equation}
		\delta m^2_{\mathrm{YUCT}} = \frac{\lambda}{32\pi^2}\,\LUV^2 \int_0^{\infty} e^{-2t^{2/3}} dt + \mathcal{O}(m^2/\LUV^2).
		\label{eq:scalar_lead}
	\end{equation}
	
	先前（版本 3.1）报告积分 \(C_0 \equiv \int_0^{\infty} e^{-2t^{2/3}} dt \approx 0.626\)。这个值是不正确的。我们现在精确计算它。
	
	\subsection{\(C_0\) 的精确计算}
	
	\begin{theorem}[\(C_0\) 的精确值]
		\label{thm:C0}
		\[
		C_0 = \int_0^{\infty} e^{-2t^{2/3}} dt = \frac{3\sqrt{\pi}}{8\sqrt{2}} \approx 0.46999.
		\]
	\end{theorem}
	
	\begin{proof}
		令 \(u = 2t^{2/3}\)。则 \(t = (u/2)^{3/2}\)，且
		\[
		dt = \frac{3}{2} (u/2)^{1/2} \cdot \frac{1}{2} du = \frac{3}{4\sqrt{2}} u^{1/2} du.
		\]
		因此
		\[
		C_0 = \int_0^{\infty} e^{-u} \cdot \frac{3}{4\sqrt{2}} u^{1/2} du = \frac{3}{4\sqrt{2}} \int_0^{\infty} u^{1/2} e^{-u} du = \frac{3}{4\sqrt{2}} \, \Gamma(3/2).
		\]
		利用 \(\Gamma(3/2) = \frac{\sqrt{\pi}}{2}\)，得到
		\[
		C_0 = \frac{3}{4\sqrt{2}} \cdot \frac{\sqrt{\pi}}{2} = \frac{3\sqrt{\pi}}{8\sqrt{2}} \approx 0.46999.
		\]
	\end{proof}
	
	\subsection{修正后的标量质量修正}
	
	采用正确的 \(C_0\) 值，标量质量修正的主导项为
	\begin{equation}
		\boxed{\delta m^2_{\mathrm{YUCT}} = \frac{\lambda}{32\pi^2}\,\LUV^2 \cdot \frac{3\sqrt{\pi}}{8\sqrt{2}} + \mathcal{O}(m^2/\LUV^2)}.
		\label{eq:scalar_corrected}
	\end{equation}
	
	对于希格斯玻色子，\(\lambda = m_H^2/(2v^2) \approx 0.13\)（其中 \(v \approx 246\) GeV）。采用普朗克尺度等同 \(\LUV = 3\pi\Mpl \approx 1.15\times 10^{20}\) GeV：
	\begin{align}
		\delta m^2_H &\approx \frac{0.13}{32\pi^2} \cdot (1.15\times 10^{20})^2 \cdot 0.46999 \notag\\
		&\approx 2.47 \times 10^{36}\ \mathrm{GeV}^2, \notag\\
		\delta m_H &\approx 1.57 \times 10^{18}\ \mathrm{GeV}.
		\label{eq:planck_correction}
	\end{align}
	
	这个值超过观测到的希格斯质量（\(m_H^{\mathrm{exp}} \approx 125\) GeV）达 \(1.26 \times 10^{16}\) 倍。如果拉格朗日量中的裸质量为零，则物理质量完全由这个辐射修正决定，预测值显然与观测不符。
	
	% ============================================================
	\section{物理含义：两种情景}
	\label{sec:scenarios}
	
	修正后的计算迫使我们直面等级问题。我们考察两种逻辑上自洽的替代方案。
	
	\subsection{情景 A：具有有限反项的普朗克尺度网络}
	
	在此情景中，YUCT 保留等同 \(\LUV = 3\pi\Mpl\)。任何粒子的物理质量由 YUCT 质量公式给出
	\begin{equation}
		m_{\mathrm{phys}} = \Mpl \left(\frac{2}{3}\right)^L \xi_f,
		\label{eq:mass_formula}
	\end{equation}
	其中 \(L\) 是协调深度，\(\xi_f\) 是共振因子（见附录 J）。拉格朗日量中的裸质量是一个自由参数，被选择来抵消辐射修正并重现观测质量。
	
	从数学上讲，这是完全自洽的：理论是有限的，所需的反项也是有限的，从未遇到无穷大。YUCT 质量公式 Eq.~(\ref{eq:mass_formula}) 提供了固定反项的重整化条件。
	
	然而，从物理上讲，此情景要求裸质量与辐射修正以 \(10^{16}\) 分之一的高精度相互抵消——这正是等级问题试图避免的那种精细调节。YUCT 没有解释为何要选择裸质量来实现这种抵消；它只是用巨大的有限减法替换了通常的无穷大减法。
	
	\subsection{情景 B：标度不变的起源与 TeV 尺度的网络}
	
	另一种在审美上更具吸引力的解决方式是放弃 \(\LUV\) 等于普朗克尺度的等同。如果普朗克尺度的基本理论不包含裸质量项（即它是标度不变的），那么所有质量都是通过协调形状因子辐射产生的。在这种情况下，物理的希格斯质量\textit{就是}辐射修正：
	\begin{equation}
		m_H^2 = \delta m^2_{\mathrm{YUCT}} = \frac{\lambda}{32\pi^2}\,\LUV^2 \cdot C_0.
		\label{eq:scale_inv}
	\end{equation}
	
	解出 \(\LUV\)：
	\begin{equation}
		\LUV = \sqrt{\frac{32\pi^2\, m_H^2}{\lambda\, C_0}}.
		\label{eq:LUV_solve}
	\end{equation}
	
	代入修正后的 \(C_0 = 3\sqrt{\pi}/(8\sqrt{2}) \approx 0.46999\)，\(\lambda = 0.13\)，以及 \(m_H = 125\) GeV，得到
	\begin{equation}
		\boxed{\LUV \approx 8.96\ \mathrm{TeV}}.
		\label{eq:LUV_TeV}
	\end{equation}
	
	这个值是理论的预测，而不是自由参数。它将基本协调尺度置于 TeV 范围，使其在实验上可及。
	
	\subsection{情景 B 的物理解释}
	
	如果 \(\LUV \approx 9\) TeV，协调网络在电弱尺度处“结晶”，而不是在普朗克尺度。形状因子 \(F(q^2) = \exp[-(q^2/\LUV^2)^{2/3}]\) 在几个 TeV 以上开始压制虚动量。这具有直接的可观测后果：
	
	\begin{enumerate}[label=(\arabic*)]
		\item \textbf{LHC 上的 Drell-Yan 和双喷注截面} 与标准模型预测相比，在高不变质量区应表现出事件缺失。对于 \(m_{jj} \gtrsim 5\) TeV，压制变得显著。
		\item \textbf{希格斯对产生} 和矢量玻色子散射应表现出反常的能量依赖，因为形状因子修改了希格斯自耦合。
		\item \textbf{电弱精确观测量中的偏差} 可能被预期，但由于 LEP 测量所探测的尺度远低于 \(\LUV\)，这些偏差可能很小。
	\end{enumerate}
	
	关键是，此情景中的最小长度为 \(\Lmin = 2\pi\hbar c/\LUV \approx 1.38\times 10^{-19}\) m，这比普朗克长度（\(\sim 10^{-35}\) m）大得多。这意味着时空的颗粒度不是一个量子引力效应，而是一个电弱尺度现象，可能与希格斯机制本身相关。
	
	\subsection{电弱尺度与宇宙学常数之间的分形桥梁}
	\label{subsec:bridge}
	
	如果情景 B 正确且基本协调尺度为 \(\LUV \approx 8.96\) TeV，那么它与宇宙学之间就产生了深刻的联系。在 YUCT 中，真空不是单一的连续介质，而是一个分形的协调层等级结构。尺度 \(\LUV\) 刻画了标准模型场经历协调形状因子的那个层。然而，还存在更深的（红外）层，对应于逐步降低的能量尺度。宇宙学常数，即观测到的暗能量密度，由最深层的红外层产生，其尺度 \(\Lambda_{\mathrm{IR}}\) 由以下要求决定：所有量子场的零点能对整个分形等级结构积分后，再现观测值 \(\rho_{\Lambda} \approx 2.51 \times 10^{-47}\) GeV\(^4\)（见附录 M）。
	
	UV 尺度 \(\LUV\) 与 IR 尺度 \(\Lambda_{\mathrm{IR}}\) 之间的关系由 YUCT 的基本标度量子决定：
	\begin{equation}
		q = \left(\frac{3}{2}\right)^{1/3} \approx 1.1447.
	\end{equation}
	这两个尺度由整数个 \(N\) 协调层分隔：
	\begin{equation}
		\LUV = \Lambda_{\mathrm{IR}} \cdot q^{N}.
		\label{eq:bridge}
	\end{equation}
	
	从观测到的宇宙学常数可提取出 \(\Lambda_{\mathrm{IR}} \approx 0.0032\) eV（见附录 M，Eq.~(M.47)）。取 \(\LUV \approx 8.96 \times 10^{12}\) eV，得到
	\begin{equation}
		N = \frac{\ln(\LUV / \Lambda_{\mathrm{IR}})}{\ln q}
		= \frac{\ln(8.96 \times 10^{12} / 0.0032)}{\frac{1}{3}\ln(3/2)}
		\approx \frac{35.564}{0.135155} \approx 263.1.
		\label{eq:N_layers}
	\end{equation}
	
	四舍五入到最接近的半整数（如 YUCT 的代数循环对费米子层的要求），我们得到基本结构数
	\begin{equation}
		\boxed{N = 263.0}.
	\end{equation}
	
	这是一个非凡的结果。电弱尺度（希格斯玻色子的质量）与宇宙学常数恰好通过 \(263\) 个离散的协调等级步骤联系起来，没有任何自由参数。宇宙学常数的观测微小值不是一个谜团；它只是反映了将粒子物理学与宇宙学分开的协调网络的巨大深度。
	
	\textbf{解释：} 等级结构中的每一层都充当一个“尺度变换器”：给定层的零点能被下一个更深层的协调形状因子压制。总真空能是所有层贡献之和，每一层比前一层小 \(\beta = 2/3\) 倍。几何级数收敛到一个有限值，正是观测到的暗能量密度（见附录 M）。整数 \(N \approx 263\) 是位于电弱尺度与中微子质量尺度（最深可及物理层）之间的层数。这个数恰好是半整数是 YUCT 框架的一个非平凡一致性检验，进一步支持了情景 B。
	
	\subsection{两种情景的比较}
	
	\begin{table}[H]
		\centering
		\caption{情景 A 与情景 B 的比较}
		\label{tab:scenarios}
		\small
		\begin{tabular}{@{}lcc@{}}
			\toprule
			\textbf{性质} & \textbf{情景 A} & \textbf{情景 B} \\
			\midrule
			UV 尺度 \(\LUV\) & \(3\pi\Mpl \approx 10^{20}\) GeV & \(\approx 8.96\) TeV \\
			最小长度 \(\Lmin\) & \(\frac{2}{3}\Lpl \approx 10^{-35}\) m & \(\approx 10^{-19}\) m \\
			拉格朗日量中的裸质量 & 非零（精细调节） & 零（标度不变性） \\
			辐射修正 \(\delta m_H\) & \(10^{18}\) GeV & 125 GeV \\
			实验可及性 & 普朗克尺度（不可及） & LHC 能量（可及） \\
			需要精细调节 & 是（\(10^{16}\) 分之一） & 否 \\
			\bottomrule
		\end{tabular}
	\end{table}
	
	% ============================================================
	\section{情景 B 的可证伪预测}
	\label{sec:predictions}
	
	如果情景 B 正确，以下定量预测可以在 LHC 和未来的对撞机上得到检验：
	
	\begin{enumerate}[label=(\arabic*)]
		\item \textbf{Drell-Yan 截面的压制。} 测量的 Drell-Yan 截面与标准模型截面之比作为双轻子不变质量 \(m_{\ell\ell}\) 的函数应遵循
		\[
		R(m_{\ell\ell}) \approx \exp\!\Bigl[-2\Bigl(\frac{m_{\ell\ell}^2}{\LUV^2}\Bigr)^{\!2/3}\,\Bigr],
		\]
		其中 \(\LUV \approx 8.96\) TeV。在 \(m_{\ell\ell} \approx 5.4\) TeV 处预测有 10\% 的压制，在 \(m_{\ell\ell} \approx 10\) TeV 处增长到约 2 倍。
		
		\item \textbf{希格斯对产生。} 形状因子在高动量转移下修改希格斯自耦合。双希格斯不变质量分布应在 \(m_{HH} \gtrsim 3\) TeV 时偏离标准模型。
		
		\item \textbf{费米子质量。} 如果所有费米子质量也是辐射产生的，则它们的汤川耦合必须被选择使得辐射修正重现观测到的质量谱。这施加了一个自洽条件，将费米子协调深度 \(L_f\) 和共振因子 \(\xi_f\) 与辐射修正联系起来。详细分析留待未来工作。
	\end{enumerate}
	
	\subsection{关于精细结构常数的一个注记}
	
	从拓扑不变量 \(N_{\mathrm{gauge}} = 12\)（附录 G）推导 \(\alpha^{-1} \approx 137.036\) 与 \(\LUV\) 的值无关。规范耦合常数是对数跑动的；它们在 \(\LUV\) 处的初值由 \(F_{18}\) 的拓扑固定。在情景 A 中，\(\LUV \sim \Mpl\)，从普朗克尺度到电弱尺度的对数跑动给出 \(\alpha^{-1} \approx 137.036\)。在情景 B 中，\(\LUV \sim 9\) TeV，跑动区间短得多。那么对 \(\alpha^{-1}\) 的拓扑预测将在 TeV 尺度适用，而低能值将通过小的对数演化得到。这将改变 \(\alpha^{-1}\) 的数值预测，必须与 CODATA 值核对。初步估计表明偏移在拓扑预测的当前不确定度之内，但精确计算是必要的，并将在附录 G 的未来修订版中给出。
	
	% ============================================================
	\section{讨论：YUCT 是万有理论还是电弱尺度理论？}
	\label{sec:discussion}
	
	最初的 YUCT 框架试图通过单一的协调等级结构解释从普朗克质量到电子质量的所有物理尺度。修正后的标量质量计算揭示了这一雄心中的一个矛盾：如果协调网络在普朗克尺度上运作，辐射修正是巨大的，必须通过精细调节的裸质量来抵消。
	
	情景 B 通过重新解释 YUCT 质量阶梯解决了这一矛盾。阶梯不是直接通过深度 \(L\) 产生质量，而是决定了\textbf{汤川耦合}，后者反过来在尺度 \(\LUV \sim 9\) TeV 上辐射产生质量。普朗克尺度于是成为协调网络\textit{开始}形成的尺度，但物理的颗粒度——形状因子变得重要的尺度——要低得多。
	
	这种重新解释保留了 YUCT 的代数之美，同时消除了精细调节问题。它还将 YUCT 从一个思辨性的量子引力理论转变为一个可证伪的电弱物理模型，可以在当前和未来的对撞机上检验。
	
	% ============================================================
	\section{结论}
	\label{sec:conclusion}
	
	我们修正了 YUCT 内标量质量修正中的一个计算错误，得到了精确的解析值 \(C_0 = 3\sqrt{\pi}/(8\sqrt{2}) \approx 0.46999\)。采用普朗克尺度等同 \(\LUV = 3\pi\Mpl\)，希格斯质量的辐射修正为 \(\delta m_H \approx 1.57\times 10^{18}\) GeV，需要与裸质量进行精细调节的抵消。
	
	作为替代方案，我们提出了情景 B，其中普朗克尺度下的裸拉格朗日量是标度不变的，所有质量都是辐射产生的。自洽性要求基本协调尺度为 \(\LUV \approx 8.96\) TeV——一个 LHC 可及的值。这一情景对高能散射过程给出了具体的、可证伪的预测。
	
	在情景 A（审美上不令人满意，但与 YUCT 的普朗克尺度起源一致）和情景 B（具有预测性且可检验，但需要重新解释 YUCT 质量阶梯）之间的选择不能仅靠理论依据做出。它必须由实验决定。我们敦促 LHC 合作组检查高不变质量的 Drell-Yan 和双喷注数据，寻找 Eq.~(\ref{eq:F}) 在 \(\LUV \sim 9\) TeV 时预测的形状因子压制的证据。零结果将排除情景 B，并且尽管情景 A 在逻辑上仍然可能，但这将削弱 YUCT 作为等级问题解决方案的说服力。
	
	\section*{致谢}
	作者感谢 YUCT 研讨会参与者的批判性讨论。
	
	\section*{数据可用性}
	所有计算都是自洽的；用于数值评估的 Python 代码可在 YPSDC 存储库中获得：\url{https://github.com/Alexey-Yakushev-YUCT/YPSDC}。
	
	\begin{thebibliography}{99}
		\bibitem{AppendixL} A.\,V.\,Yakushev, “附录 L：分形协调误差标度,” 载于 \emph{YUCT}, Zenodo (2026).
		\bibitem{AppendixV} A.\,V.\,Yakushev, “附录 V：时间作为协调过程的熵,” 载于 \emph{YUCT}, Zenodo (2026).
		\bibitem{AppendixG} A.\,V.\,Yakushev, “附录 G：作为协调网络几何张力的引力与 YUCT 中的循环宇宙学,” 载于 \emph{YUCT}, Zenodo (2026).
		\bibitem{AppendixM} A.\,V.\,Yakushev, “附录 M：YUCT 宇宙学——作为协调相变的暴胀,” 载于 \emph{YUCT}, Zenodo (2026).
		\bibitem{AppendixLambda} A.\,V.\,Yakushev, “附录 \(\Lambda\)：YUCT 中等级问题和宇宙学常数问题的解决,” 载于 \emph{YUCT}, Zenodo (2026).
		\bibitem{AppendixJ} A.\,V.\,Yakushev, “附录 J：从 YUCT 拓扑偏移完整推导标准模型味参数,” 载于 \emph{YUCT}, Zenodo (2026).
		\bibitem{AppendixFerra} A.\,V.\,Yakushev, “附录 Ferra1.2：有序相和无序相的普适协调常数,” 载于 \emph{YUCT}, Zenodo (2026).
	\end{thebibliography}
	
\end{document}